\documentclass{ca-capstone}
\capstonerefs{crandall-pomerance,hps,msu-ent,sagemath}
\begin{document}
% Generated by `make bench CAPSTONE=primality-testing` (taskgen.capstone report) from
% bench/results.json: reference time x slack on the reference machine.
% Do not edit: change bench/config.yaml and re-run the benchmark.
\newcommand{\ReqSlack}{10}
\newcommand{\ReqBitsMax}{128}
\newcommand{\ReqSeconds}{0{,}01}

\capstonetitle{Тесты простоты и псевдопростые числа}

\goal
Требуется реализовать вероятностные тесты простоты, исследовать их свидетелей и
классы псевдопростых чисел.

\functionality
\begin{enumerate}
  \item Тест Вильсона как эталон корректности на малых входах, без требования к скорости.
  \item Тесты Ферма и Миллера-Рабина.
  \item Тест Соловея-Штрассена с самостоятельно реализованным символом Якоби.
  \item Критерий Корсельта для чисел Кармайкла.
  \item Подсчёт оснований Ферма-лжецов для свободного от квадратов~$N$ по заданной
    факторизации.
  \item Тест Люка-Лемера с быстрым приведением по модулю числа Мерсенна, без тривиального
    деления на это число.
\end{enumerate}

\requirements
\begin{enumerate}
  \item Прохождение вероятностного теста должно трактоваться как «вероятно простое», а не как
    доказательство простоты. Каждый свидетель непростоты проверяется независимо.
  \item Реализация должна корректно обрабатывать малые простые и составные числа, чётные входные числа, $N<2$ и все перечисленные семейства псевдопростых; 
    демонстрировать повторяемость результатов при фиксированных основаниях.
  \item Смешанные тестовые запросы с кандидатами длиной до \textbf{\ReqBitsMax~бит} должны
    выполняться не более чем за \textbf{\ReqSeconds~с}; тест Вильсона проводится только на
    малых входах.
\end{enumerate}

\details
Разрешается использовать любой язык программирования, в частности Python, систему 
компьютерной алгебры SageMath, встроенные большие целые и вспомогательное модульное возведение
в степень. Библиотечные \texttt{is\_prime}, символ Якоби и тест Люка-Лемера в решении
запрещены и используются как оракул. Факторизация для критерия Корсельта и подсчёта лжецов
в работе не выполняется: разложение включается в набор входных данных.

\schedule
Работа выполняется индивидуально. Срок выполнения -- \textbf{месяц}; ожидаемая трудоёмкость
-- примерно 38-42 часов. Реализация должна быть разработана и написана вами.

\remarks
\begin{enumerate}
  \item Для каждого вероятностного теста необходимо показать как принятые, так и отвергнутые основания.
  \item Для свободного от квадратов~$N$ проверяется формула числа лжецов
    $\prod_{p\mid N}\gcd(N-1,p-1)$.
  \item Отдельно тестируются приведение Мерсенна и финальная невязка Люка-Лемера.
\end{enumerate}

\references
\end{document}
